\documentclass[10pt,a4paper]{amsart}
\usepackage{amsmath,amssymb,amsthm,graphicx}
% Retained preamble for the Stacks Project chapter snapshot in this directory.
% The chapter's own preamble.tex is a copy of the shared Stacks Project preamble; the block
% below is sliced MECHANICALLY, line for line, from it. Nothing was added, removed or
% redefined; the wrapper class, the xr-hyper/hyperref wiring and the externaldocument list
% are omitted because the assembler supplies the document class.
%
% lines taken from preamble.tex: 3, 5, 9, 11, 12, 156, 157, 158, ... (51 lines total)
%
% package coverage restored from the chapter's own preamble (2026-10-13): xy, xy-2cell

\usepackage[all]{xy}
\xyoption{2cell}
\UseAllTwocells

\newenvironment{reference}{\comment}{\endcomment}
\newenvironment{slogan}{\comment}{\endcomment}
\newenvironment{history}{\comment}{\endcomment}
\theoremstyle{plain}
\newtheorem{theorem}[subsection]{Theorem}
\newtheorem{proposition}[subsection]{Proposition}
\newtheorem{lemma}[subsection]{Lemma}
\theoremstyle{definition}
\newtheorem{definition}[subsection]{Definition}
\newtheorem{example}[subsection]{Example}
\newtheorem{exercise}[subsection]{Exercise}
\newtheorem{situation}[subsection]{Situation}
\theoremstyle{remark}
\newtheorem{remark}[subsection]{Remark}
\newtheorem{remarks}[subsection]{Remarks}
\numberwithin{equation}{subsection}
\def\lim{\mathop{\mathrm{lim}}\nolimits}
\def\colim{\mathop{\mathrm{colim}}\nolimits}
\def\Spec{\mathop{\mathrm{Spec}}}
\def\Hom{\mathop{\mathrm{Hom}}\nolimits}
\def\Ext{\mathop{\mathrm{Ext}}\nolimits}
\def\SheafHom{\mathop{\mathcal{H}\!\mathit{om}}\nolimits}
\def\SheafExt{\mathop{\mathcal{E}\!\mathit{xt}}\nolimits}
\def\Sch{\mathit{Sch}}
\def\Mor{\mathop{\mathrm{Mor}}\nolimits}
\def\Ob{\mathop{\mathrm{Ob}}\nolimits}
\def\Sh{\mathop{\mathit{Sh}}\nolimits}
\def\NL{\mathop{N\!L}\nolimits}
\def\CH{\mathop{\mathrm{CH}}\nolimits}
\def\proetale{{pro\text{-}\acute{e}tale}}
\def\etale{{\acute{e}tale}}
\def\QCoh{\mathit{QCoh}}
\def\Ker{\mathop{\mathrm{Ker}}}
\def\Im{\mathop{\mathrm{Im}}}
\def\Coker{\mathop{\mathrm{Coker}}}
\def\Coim{\mathop{\mathrm{Coim}}}
\DeclareMathSymbol{\boxtimes}{\mathbin}{AMSa}{"02}
\def\QCohstack{\mathcal{QC}\!\mathit{oh}}
\def\Cohstack{\mathcal{C}\!\mathit{oh}}
\def\Spacesstack{\mathcal{S}\!\mathit{paces}}
\def\Quotfunctor{\mathrm{Quot}}
\def\Hilbfunctor{\mathrm{Hilb}}
\def\Curvesstack{\mathcal{C}\!\mathit{urves}}
\def\Polarizedstack{\mathcal{P}\!\mathit{olarized}}
\def\Complexesstack{\mathcal{C}\!\mathit{omplexes}}
\def\Pic{\mathop{\mathrm{Pic}}\nolimits}
\def\Picardstack{\mathcal{P}\!\mathit{ic}}
\def\Picardfunctor{\mathrm{Pic}}
\def\Deformationcategory{\mathcal{D}\!\mathit{ef}}
\usepackage{cleveref}
\crefname{originallemma}{original lemma}{original lemmas}
\Crefname{originallemma}{Original lemma}{Original lemmas}
\makeatletter
\newlabel{derived-lemma-fully-faithful-adjoint-kernel-zero}{{09J1}{source}{}{sourcepointer-1}{}}
\newlabel{derived-lemma-fully-faithful-adjoint-kernel-zero@cref}{{[originallemma][0][]09J1}{[1][0][]source}{}{}{}}
\makeatother
\crefname{originallemma}{original lemma}{original lemmas}
\Crefname{originallemma}{Original lemma}{Original lemmas}
\makeatletter
\newlabel{lemma-formal-glueing-on-closed}{{0AEU}{source}{}{sourcepointer-2}{}}
\newlabel{lemma-formal-glueing-on-closed@cref}{{[originallemma][0][]0AEU}{[1][0][]source}{}{}{}}
\makeatother
\crefname{originallemma}{original lemma}{original lemmas}
\Crefname{originallemma}{Original lemma}{Original lemmas}
\makeatletter
\newlabel{situation-formal-glueing}{{0AEV}{source}{}{sourcepointer-3}{}}
\newlabel{situation-formal-glueing@cref}{{[originallemma][0][]0AEV}{[1][0][]source}{}{}{}}
\makeatother
\crefname{originallemma}{original lemma}{original lemmas}
\Crefname{originallemma}{Original lemma}{Original lemmas}
\makeatletter
\newlabel{spaces-perfect-lemma-quasi-coherence-direct-image}{{08FA}{source}{}{sourcepointer-4}{}}
\newlabel{spaces-perfect-lemma-quasi-coherence-direct-image@cref}{{[originallemma][0][]08FA}{[1][0][]source}{}{}{}}
\makeatother
\crefname{originallemma}{original lemma}{original lemmas}
\Crefname{originallemma}{Original lemma}{Original lemmas}
\makeatletter
\newlabel{spaces-perfect-lemma-quasi-coherence-pullback}{{08F4}{source}{}{sourcepointer-5}{}}
\newlabel{spaces-perfect-lemma-quasi-coherence-pullback@cref}{{[originallemma][0][]08F4}{[1][0][]source}{}{}{}}
\makeatother
\title{Rf\_* is an equivalence between supported derived categories with quasi-inverse Lf\textasciicircum{}*}
\author{The Stacks Project Authors}
\date{stacks-project:spaces-pushouts.tex:local-snapshot-20261012; GNU FDL 1.2 or later}
\begin{document}
\maketitle
\noindent\textit{Editorial note.} Statement span is lines 2243-2249 of sources/stacks-spaces-pushouts/pinned-original-source.tex; proof span is lines 2251-2269, the complete author proof (17 body lines).\par\medskip
\noindent\textit{Editorial note.} Navigation stubs: 5 cited-but-not-included labels are supplied as non-load-bearing pointers with their pinned tags where the snapshot tag index has them.\par\medskip
\noindent\textit{Editorial note.} extra\_packages is empty on purpose: the retained preamble already loads xy with options and a second declaration clashes in the assembler wrapper.\par\medskip
\noindent\textit{Editorial note.} Text is retained verbatim, including the author's own typos and placeholders.\par\medskip
\noindent\textit{Editorial note.} Queue index 31 of queue\_spaces-pushouts.json (pb=921, stmt 2243-2249, proof 2251-2269).\par\medskip
\noindent\textit{Editorial note.} spaces-pushouts wave-1: fresh adjudication accept (2026-10-14); six-check dispatched pool (§2 row appended 02:4x).\par\medskip
\noindent\textit{Allowed standard prerequisites (editorial).}\begin{itemize}
\item Explicit prerequisite boundary. The cited inputs are prerequisites; the argument in the retained span is not replaced by any cited result.
\end{itemize}\medskip
\ifdefined\hypertarget\hypertarget{sourcepointer-1}{}\fi
\noindent\textit{Source reference.} Original lemma 09J1: sources/stacks-spaces-pushouts/pinned-original-source.tex, source lines 2269--2269. Cited by the retained statement or proof but not part of the collected spans; supplied only as a navigation pointer. The argument does not depend on its content.\par\medskip
\ifdefined\hypertarget\hypertarget{sourcepointer-2}{}\fi
\noindent\textit{Source reference.} Original lemma 0AEU: sources/stacks-spaces-pushouts/pinned-original-source.tex, source lines 1994--2000. Cited by the retained statement or proof but not part of the collected spans; supplied only as a navigation pointer. The argument does not depend on its content.\par\medskip
\ifdefined\hypertarget\hypertarget{sourcepointer-3}{}\fi
\noindent\textit{Source reference.} Original lemma 0AEV: sources/stacks-spaces-pushouts/pinned-original-source.tex, source lines 2046--2060. Cited by the retained statement or proof but not part of the collected spans; supplied only as a navigation pointer. The argument does not depend on its content.\par\medskip
\ifdefined\hypertarget\hypertarget{sourcepointer-4}{}\fi
\noindent\textit{Source reference.} Original lemma 08FA: sources/stacks-spaces-pushouts/pinned-original-source.tex, source lines 2269--2269. Cited by the retained statement or proof but not part of the collected spans; supplied only as a navigation pointer. The argument does not depend on its content.\par\medskip
\ifdefined\hypertarget\hypertarget{sourcepointer-5}{}\fi
\noindent\textit{Source reference.} Original lemma 08F4: sources/stacks-spaces-pushouts/pinned-original-source.tex, source lines 2269--2269. Cited by the retained statement or proof but not part of the collected spans; supplied only as a navigation pointer. The argument does not depend on its content.\par\medskip
\begin{lemma}
\label{lemma-derived-equivalent}
In Situation \ref{situation-formal-glueing} the functor
$Rf_*$ induces an equivalence between $D_{\QCoh, |f^{-1}Z|}(\mathcal{O}_Y)$
and $D_{\QCoh, |Z|}(\mathcal{O}_X)$ with quasi-inverse given by
$Lf^*$.
\end{lemma}

\begin{proof}
Since $f$ is quasi-compact and quasi-separated we see that $Rf_*$
defines a functor from $D_{\QCoh, |f^{-1}Z|}(\mathcal{O}_Y)$
to $D_{\QCoh, |Z|}(\mathcal{O}_X)$, see
Derived Categories of Spaces, Lemma
\ref{spaces-perfect-lemma-quasi-coherence-direct-image}.
By Derived Categories of Spaces, Lemma
\ref{spaces-perfect-lemma-quasi-coherence-pullback}
we see that $Lf^*$ maps $D_{\QCoh, |Z|}(\mathcal{O}_X)$
into $D_{\QCoh, |f^{-1}Z|}(\mathcal{O}_Y)$.
In Lemma \ref{lemma-formal-glueing-on-closed} we have seen that
$Lf^*Rf_*Q = Q$ for $Q$ in $D_{\QCoh, |f^{-1}Z|}(\mathcal{O}_Y)$.
By the dual of Derived Categories, Lemma
\ref{derived-lemma-fully-faithful-adjoint-kernel-zero}
to finish the proof it suffices to show that $Lf^*K = 0$
implies $K = 0$ for $K$ in $D_{\QCoh, |Z|}(\mathcal{O}_X)$.
This follows from the fact that $f$ is flat at all points of
$f^{-1}Z$ and the fact that $f^{-1}Z \to Z$ is surjective.
\end{proof}

\end{document}
